% ============================================================
%  Book chapter 17: Text Data
% ============================================================
\chapter{Text Data}

\textbf{What this chapter gives you:} a clear picture of what a Python \code{str} really is
(a sequence of Unicode code points), how text turns into bytes and back, and how to pull fields
out of messy text with regular expressions. You will need both skills every time you read model
output, log files, scraped pages or JSON from an API.

% ------------------------------------------------------------
\tikzset{c17cell/.style={draw=cgray,minimum width=5.2mm,minimum height=6.5mm,inner sep=0pt,
    font=\ttfamily\small},
  c17byte/.style={draw=cblue,fill=cblue!10,minimum width=11mm,minimum height=7mm,inner sep=1pt,
    font=\ttfamily\small},
  c17chr/.style={draw=cgreen,fill=cgreen!10,minimum width=11mm,minimum height=7mm,inner sep=1pt,
    font=\ttfamily\small}}

\section{Text Strings: Unicode}

\subsection{Python Unicode strings}

A Python 3 \code{str} is a sequence of \emph{code points}. A code point is just a number that
Unicode assigns to one character. ASCII covers only 128 of them. Unicode covers more than
150{,}000: accented letters, other alphabets, math symbols, emoji.
You can write any character with an escape sequence, so your source file can stay plain ASCII.

\begin{center}
\begin{tabularx}{\linewidth}{@{}l l X@{}}
\toprule
\textbf{Escape} & \textbf{Example} & \textbf{Meaning} \\
\midrule
\code{\textbackslash u} + 4 hex digits & \code{"\textbackslash u00e9"} & code point up to \code{FFFF} \\
\code{\textbackslash U} + 8 hex digits & \code{"\textbackslash U0001f47b"} & any code point, pad with zeros \\
\code{\textbackslash N\{name\}} & \code{"\textbackslash N\{SNOWMAN\}"} & the official Unicode name (any case) \\
\bottomrule
\end{tabularx}
\end{center}

The \code{unicodedata} module converts between names and characters. The built-ins \code{ord()}
and \code{chr()} convert between a character and its number.

\begin{pycode}
import unicodedata

a = "caf\u00e9"
b = "caf\N{LATIN SMALL LETTER E WITH ACUTE}"
print(a == b, len(a))                       # same text, 4 characters
print(unicodedata.name("\u00e9"))
print(hex(ord("\u00e9")), ascii(chr(233)))
print(unicodedata.lookup("snowman") == "\u2603")
print(len("\U0001f47b"))                    # one emoji is one character
\end{pycode}
\begin{rawcode}
True 4
LATIN SMALL LETTER E WITH ACUTE
0xe9 '\xe9'
True
1
\end{rawcode}

\begin{typecheck}
\code{len(s)} counts \emph{code points}, not bytes and not what you see on screen. One emoji is
length 1. A \code{str} has no encoding at all. It is only when you turn it into \code{bytes} that
an encoding appears.
\end{typecheck}

\begin{trap}
Your terminal or font may not be able to \emph{show} a character that Python holds correctly.
A box or question mark on screen does not mean the string is broken. Use \code{ascii(s)} or
\code{s.encode("unicode-escape")} to see what is really inside. On Windows, \code{print()} of
non-ASCII text can even raise an error if the console uses an old code page.
\end{trap}

\subsection{UTF-8: encode and decode}

Files, sockets and APIs carry \code{bytes}, not \code{str}. Two operations cross the border:

\begin{center}
\begin{tikzpicture}[box/.style={draw=#1,fill=#1!10,rounded corners=3pt,minimum width=26mm,
    minimum height=9mm,font=\ttfamily\small},>=Stealth]
  \node[box=cgreen] (s) at (0,0) {str};
  \node[box=cblue]  (b) at (7,0) {bytes};
  \draw[->,thick] (s.east) -- node[above,font=\ttfamily\footnotesize]{.encode("utf-8")} (b.west);
  \draw[->,thick] ([yshift=-3mm]b.west) -- node[below,font=\ttfamily\footnotesize]{.decode("utf-8")} ([yshift=-3mm]s.east);
\end{tikzpicture}
\end{center}
% MD-BEGIN
% ```mermaid
% flowchart LR
%     S["str (code points)"] -- ".encode('utf-8')" --> B["bytes (0-255 values)"]
%     B -- ".decode('utf-8')" --> S
% ```
% MD-END

UTF-8 uses one to four bytes per character. ASCII characters take one byte, so plain English
text looks the same in ASCII and UTF-8. Accented Latin letters take two, most other scripts
three, emoji four. It is the default almost everywhere, so use it unless you are forced not to.

The picture below shows what happens to \code{"caf\textbackslash u00e9"}: four characters become five bytes,
because the last character needs two.

\begin{center}
\begin{tikzpicture}[>=Stealth]
  \node[font=\sffamily\footnotesize,anchor=east] at (-0.5,1.8) {str: 4 characters};
  \node[c17chr] (c1) at (0,1.8) {c};
  \node[c17chr] (c2) at (1.3,1.8) {a};
  \node[c17chr] (c3) at (2.6,1.8) {f};
  \node[c17chr,minimum width=26mm] (c4) at (4.55,1.8) {\'e\ \ U+00E9};
  \node[font=\sffamily\footnotesize,anchor=east] at (-0.5,0) {bytes: 5 values};
  \node[c17byte] (b1) at (0,0) {63};
  \node[c17byte] (b2) at (1.3,0) {61};
  \node[c17byte] (b3) at (2.6,0) {66};
  \node[c17byte] (b4) at (3.9,0) {c3};
  \node[c17byte] (b5) at (5.2,0) {a9};
  \foreach \i in {1,2,3} \draw[->,thick,cgray] (c\i.south) -- (b\i.north);
  \draw[->,thick,cgray] (c4.south) -- (b4.north);
  \draw[->,thick,cgray] (c4.south) -- (b5.north);
  \node[font=\sffamily\scriptsize,text=cgray,anchor=west] at (6.9,1.8) {1 byte each for ASCII};
  \node[font=\sffamily\scriptsize,text=cgray,anchor=west] at (6.9,0.9) {2 bytes for \'e};
\end{tikzpicture}
\end{center}
% MD-BEGIN
% ```text
% str    :  c    a    f    é (U+00E9)
%           |    |    |     |     \
% bytes  : 63   61   66    c3     a9      (hex)   -> 4 characters, 5 bytes
% ```
% MD-END

\begin{pycode}
s = "na\u00efve \u2603"          # naive with a diaeresis, plus a snowman
data = s.encode("utf-8")
print(data, len(s), len(data))
print(type(s), type(data))
print(data.decode("utf-8") == s)
\end{pycode}
\begin{rawcode}
b'na\xc3\xafve \xe2\x98\x83' 7 10
<class 'str'> <class 'bytes'>
True
\end{rawcode}

\subsubsection*{Encode: what if the character does not fit?}

Encodings other than UTF-8 cannot hold every character. The default error handler is
\code{"strict"}, which raises \code{UnicodeEncodeError}. The second argument of \code{encode()}
picks another policy.

\begin{pycode}
for enc in ("ascii", "latin-1"):
    try:
        print(enc, s.encode(enc))
    except UnicodeEncodeError as e:
        print(enc, "failed:", e.reason)
print(s.encode("ascii", "ignore"))
print(s.encode("ascii", "replace"))
print(s.encode("ascii", "backslashreplace"))
print(s.encode("ascii", "xmlcharrefreplace"))
\end{pycode}
\begin{rawcode}
ascii failed: ordinal not in range(128)
latin-1 failed: ordinal not in range(256)
b'nave '
b'na?ve ?'
b'na\\xefve \\u2603'
b'na&#239;ve &#9731;'
\end{rawcode}

\begin{center}
\begin{tabularx}{\linewidth}{@{}l X@{}}
\toprule
\textbf{Handler} & \textbf{What it does with a character that does not fit} \\
\midrule
\code{strict} & raise an error (default) \\
\code{ignore} & drop it silently \\
\code{replace} & write \code{?} \\
\code{backslashreplace} & write a Python-style escape such as \code{\textbackslash u2603} \\
\code{xmlcharrefreplace} & write an HTML/XML numeric entity such as \code{\&\#9731;} \\
\bottomrule
\end{tabularx}
\end{center}

\subsubsection*{Decode: the encoding is a guess you must get right}

Bytes carry no label saying which encoding made them. You must know it from the source: a file
header, an HTTP \code{Content-Type}, a spec. Decode with the wrong one and you get either an
error or, worse, wrong text without any error.

\begin{pycode}
raw = "caf\u00e9".encode("utf-8")
print(raw)
print(ascii(raw.decode("latin-1")))     # no error, but wrong text
try:
    raw.decode("ascii")
except UnicodeDecodeError as e:
    print("ascii failed:", e.reason, "at", e.start)
print(ascii(b"caf\xe9".decode("utf-8", errors="replace")))
\end{pycode}
\begin{rawcode}
b'caf\xc3\xa9'
'caf\xc3\xa9'
ascii failed: ordinal not in range(128) at 3
'caf\ufffd'
\end{rawcode}

The second line is the classic garbled text (called mojibake): two strange characters where one
\code{\'e} should be. The last line shows the \code{"replace"} handler on decode: bad bytes become the
replacement character \code{U+FFFD}.

\begin{center}
\begin{tikzpicture}[>=Stealth]
  \node[c17byte] (b1) at (0,0.45) {c3};
  \node[c17byte] (b2) at (1.2,0.45) {a9};
  \node[font=\sffamily\footnotesize,anchor=east] at (-0.7,0.45) {same two bytes};
  \node[c17chr,minimum width=14mm] (u) at (6,1.2) {\'e};
  \node[c17chr,minimum width=20mm,draw=cred,fill=cred!10] (l) at (6,-0.3) {\~A\textcopyright};
  \draw[->,thick,cgreen] (1.9,0.65) -- node[above,sloped,font=\ttfamily\scriptsize]{utf-8} (u.west);
  \draw[->,thick,cred] (1.9,0.25) -- node[below,sloped,font=\ttfamily\scriptsize]{latin-1} (l.west);
  \node[font=\sffamily\scriptsize,anchor=west,text=cgreen!60!black] at (7.2,1.2) {1 character: correct};
  \node[font=\sffamily\scriptsize,anchor=west,text=cred] at (7.2,-0.3) {2 characters: mojibake};
\end{tikzpicture}
\end{center}
% MD-BEGIN
% ```mermaid
% flowchart LR
%     B["bytes c3 a9"] -- "decode utf-8" --> OK["1 character: e with acute (correct)"]
%     B -- "decode latin-1" --> BAD["2 characters: A-tilde + copyright sign (mojibake)"]
% ```
% MD-END

\begin{gamerule}[: UTF-8 everywhere]
Decode bytes to \code{str} as soon as they enter your program. Work only with \code{str}
inside. Encode to bytes as late as possible, when you write or send. Always name the encoding
explicitly (\code{"utf-8"}) instead of trusting the platform default.
\end{gamerule}

\begin{trap}
\code{errors="ignore"} and \code{errors="replace"} hide data loss. They are fine for a log
viewer. They are dangerous in a data pipeline, because corrupted training or eval text goes in
without any warning. Prefer \code{strict} and fix the cause.
\end{trap}

\begin{boss}[: the str/bytes border]
Most Unicode bugs come from one mistake: not knowing whether a value is \code{str} or
\code{bytes} at this point in the program. Keep a mental label on every variable.
Bytes arrive from outside (\code{.decode} them), text goes out (\code{.encode} it). Never
mix: \code{"a" + b"a"} raises \code{TypeError}, and \code{"a" == b"a"} is just \code{False},
not an error. If you see \code{UnicodeDecodeError}, ask first: \emph{which encoding made these
bytes?}
\end{boss}

\subsection{HTML entities}

Web text often hides characters behind entities such as \code{\&eacute;}, \code{\&\#233;} or
\code{\&\#xe9;}. The \code{html} module converts them both ways. This matters when you scrape
pages or clean model output that was copied from a web page.

\begin{pycode}
import html

print(ascii(html.unescape("&eacute; &#233; &#xe9; &amp; &lt;b&gt;")))
print(html.escape('<a href="x">Q&A</a>'))
print(ascii("caf\u00e9".encode("ascii", "xmlcharrefreplace").decode()))
\end{pycode}
\begin{rawcode}
'\xe9 \xe9 \xe9 & <b>'
&lt;a href=&quot;x&quot;&gt;Q&amp;A&lt;/a&gt;
'caf&#233;'
\end{rawcode}

\begin{typecheck}
\code{html.unescape(str)} returns \code{str}. For the other direction on a whole string, you
encode to ASCII with \code{xmlcharrefreplace} and get \code{bytes}; call \code{.decode()} to get
a \code{str} again. For the markup characters \code{< > \& "} use \code{html.escape()}.
The module \code{html.entities} also holds dictionaries of names such as \code{html5} and
\code{codepoint2name}.
\end{typecheck}

\subsection{Normalization}

Some visible characters can be built in more than one way. The letter \'e can be one code
point (\code{U+00E9}) or a plain \code{e} followed by a combining accent (\code{U+0301}). They
look identical. They are not equal.

\begin{pycode}
import unicodedata

one = "\u00e9"
two = "e\u0301"
print(one == two, len(one), len(two))
n = unicodedata.normalize("NFC", two)
print(n == one, len(n))
print(len(unicodedata.normalize("NFD", one)))
print(unicodedata.normalize("NFKC", "\ufb01"), len("\ufb01"))
\end{pycode}
\begin{rawcode}
False 1 2
True 1
2
fi 1
\end{rawcode}

\begin{center}
\begin{tikzpicture}[>=Stealth]
  \node[c17chr,minimum width=16mm] (a) at (0,0.8) {U+00E9};
  \node[c17chr] (b1) at (-0.7,-0.6) {e};
  \node[c17chr,minimum width=16mm] (b2) at (0.9,-0.6) {U+0301};
  \node[font=\sffamily\scriptsize,anchor=east] at (-1.4,0.8) {length 1};
  \node[font=\sffamily\scriptsize,anchor=east] at (-1.4,-0.6) {length 2};
  \node[draw=cteal,fill=cteal!10,rounded corners=2pt,font=\ttfamily\small] (n) at (4.6,0.1) {NFC};
  \node[c17chr,minimum width=16mm] (r) at (7.6,0.1) {U+00E9};
  \draw[->,thick,cgray] (a.east) -- (n.north west);
  \draw[->,thick,cgray] (b2.east) -- (n.south west);
  \draw[->,thick,cgray] (n) -- (r);
  \node[font=\sffamily\scriptsize,text=cred] at (0.1,-1.45) {look identical, but \code{==} is False};
  \node[font=\sffamily\scriptsize,text=cgreen!60!black] at (7.6,-0.8) {now \code{==} is True};
\end{tikzpicture}
\end{center}
% MD-BEGIN
% ```text
% "é"  = [U+00E9]            length 1  ─┐
%                                        ├─ normalize("NFC") ─> [U+00E9]  (equal)
% "é"  = [e][U+0301]         length 2  ─┘
% ```
% MD-END

\code{NFC} composes characters into the shortest form, \code{NFD} splits them apart. The
\code{NFKC} and \code{NFKD} forms also fold look-alikes (the one-character ligature \emph{fi}
becomes the two letters \code{fi}).

\begin{gamerule}
Normalize text with \code{unicodedata.normalize("NFC", s)} before you compare it, hash it, use
it as a dict key, or deduplicate a dataset. Two strings from different sources can look the same
and still be different keys.
\end{gamerule}

% ------------------------------------------------------------
\section{Text Strings: Regular Expressions}

A regular expression (regex) is a small pattern language for describing text. The standard
module is \code{re}. You give it a pattern and a source string. The pattern is itself a
\code{str}.

\subsection{The five basic functions}

\begin{center}
\begin{tabularx}{\linewidth}{@{}l l X@{}}
\toprule
\textbf{Function} & \textbf{Returns} & \textbf{Meaning} \\
\midrule
\code{re.match(p, s)} & \code{Match} or \code{None} & pattern must match at the \emph{start} of \code{s} \\
\code{re.search(p, s)} & \code{Match} or \code{None} & first match anywhere in \code{s} \\
\code{re.findall(p, s)} & \code{list} & every non-overlapping match \\
\code{re.split(p, s)} & \code{list} & pieces of \code{s} between matches \\
\code{re.sub(p, repl, s)} & \code{str} & copy of \code{s} with matches replaced \\
\bottomrule
\end{tabularx}
\end{center}

The difference between the first two functions is where they are willing to look:

\begin{center}
\begin{tikzpicture}[>=Stealth]
  \foreach \c [count=\i from 0] in {e,r,r,\textvisiblespace,4,0,4,\textvisiblespace,o,k}
    \node[c17cell] (s\i) at (\i*0.6,0) {\c};
  \node[font=\sffamily\footnotesize,anchor=east] at (-0.5,0) {source};
  \node[text=cred,font=\bfseries] at (0,-0.9) {\xmark};
  \node[font=\sffamily\scriptsize,anchor=west] at (0.5,-0.9) {\code{match(r"\textbackslash d+")}: only tries position 0, fails};
  \draw[fill=cgreen!25,draw=cgreen] (2.1,-1.45) rectangle (3.9,-1.85);
  \node[font=\sffamily\scriptsize,anchor=west] at (4.7,-1.65) {\code{search(r"\textbackslash d+")}: slides along, finds \code{404} at 4..7};
  \draw[->,thick,cgreen] (0,-1.65) -- (2.0,-1.65);
  \foreach \i in {4,5,6} \node[fill=cgreen!25,draw=cgreen,minimum width=5.2mm,minimum height=6.5mm,
     inner sep=0pt,font=\ttfamily\small] at (\i*0.6,0) {\ifnum\i=5 0\else 4\fi};
\end{tikzpicture}
\end{center}
% MD-BEGIN
% ```text
% index :  0 1 2 3 4 5 6 7 8 9
% source:  e r r _ 4 0 4 _ o k
% match : ^ tries only index 0 -> no digit there -> None
% search: ----------> slides right, first hit is "404" at span (4, 7)
% ```
% MD-END

Our sample is one log line from a training run.

\begin{pycode}
import re

line = "2025-03-14 12:01:07 INFO run=a17 model=small-7b loss=0.3125 tokens=1520"

print(re.match(r"\d{4}", line))      # starts with 4 digits
print(re.match(r"INFO", line))       # INFO is not at the start
print(re.search(r"INFO", line))      # but it is somewhere
m = re.search(r"loss=(\d+\.\d+)", line)
print(type(m))
print(m.group(), m.group(1), m.groups(), m.span())
print(float(m.group(1)) + 1)
print(re.search(r"nope", line))
if m := re.search(r"tokens=(\d+)", line):
    print(int(m.group(1)) * 2)
\end{pycode}
\begin{rawcode}
<re.Match object; span=(0, 4), match='2025'>
None
<re.Match object; span=(20, 24), match='INFO'>
<class 're.Match'>
loss=0.3125 0.3125 ('0.3125',) (48, 59)
1.3125
None
3040
\end{rawcode}

Here is what the groups of the \code{loss} pattern cover. Group 0 is the whole match, group 1
is only the part in parentheses.

\begin{center}
\begin{tikzpicture}[>=Stealth,decoration={brace,mirror,amplitude=4pt}]
  \foreach \c [count=\i from 0] in {r,u,n,\textvisiblespace,l,o,s,s,=,0,.,3,1,\textvisiblespace,o,k}
    \node[c17cell] (s\i) at (\i*0.6,0) {\c};
  \foreach \i in {0,4,9,13,16} \node[font=\sffamily\tiny,text=cgray] at (\i*0.6-0.3,0.55) {\i};
  \draw[decorate,thick,cblue] (2.1,-0.4) -- (7.5,-0.4) coordinate[pos=0.5] (g0);
  \node[font=\sffamily\scriptsize,text=cblue,below=3mm] at (g0) {group 0: \code{loss=0.31}, span (4, 13)};
  \draw[decorate,thick,cgold] (5.1,-1.5) -- (7.5,-1.5) coordinate[pos=0.5] (g1);
  \node[font=\sffamily\scriptsize,text=cgold!70!black,below=3mm] at (g1) {group 1: \code{0.31}, span (9, 13)};
\end{tikzpicture}
\end{center}
% MD-BEGIN
% ```text
% index :  0 1 2 3 4 5 6 7 8 9 10 11 12 13 14 15
% text  :  r u n _ l o s s = 0  .  3  1  _  o  k
%                  |<---- group 0 ---->|              span (4, 13)  "loss=0.31"
%                            |<- g1 ->|                span (9, 13)  "0.31"
% pattern: loss=(\d+\.\d+)
% ```
% MD-END

\begin{typecheck}
\code{match} and \code{search} return a \code{re.Match} object or \code{None}, never a plain
string. \code{m.group()} is the whole match, \code{m.group(1)} is the first parenthesised
part, \code{m.groups()} is a tuple of all of them. Every group is a \code{str}, even if it looks
like a number. Convert with \code{float()} or \code{int()}.
\end{typecheck}

\begin{trap}
Calling \code{m.group()} when \code{m} is \code{None} raises \code{AttributeError}. Always test
first (\code{if m:}) or use the walrus form shown above. And remember: \code{match} only
tries the start of the string. People expect it to behave like \code{search}.
\end{trap}

\code{findall}, \code{split} and \code{sub} work on the whole string:

\begin{pycode}
print(re.findall(r"\d+", line))
print(re.findall(r"(\w+)=(\S+)", line))
print(dict(re.findall(r"(\w+)=(\S+)", line)))
print(re.split(r"\s*,\s*", "a , b,c ,d"))
print(re.sub(r"\d", "#", "key 12345 here"))
print(re.sub(r"(\w+)=(\S+)", r"\2:\1", "a=1 b=2"))
print(re.sub(r"\d+", lambda m: str(int(m.group()) * 2), "3 apples 10 pears"))
\end{pycode}
\begin{rawcode}
['2025', '03', '14', '12', '01', '07', '17', '7', '0', '3125', '1520']
[('run', 'a17'), ('model', 'small-7b'), ('loss', '0.3125'), ('tokens', '1520')]
{'run': 'a17', 'model': 'small-7b', 'loss': '0.3125', 'tokens': '1520'}
['a', 'b', 'c', 'd']
key ##### here
1:a 2:b
6 apples 20 pears
\end{rawcode}

Notes on the output:
\begin{itemize}[nosep]
  \item The first \code{findall} grabs digits from inside \code{a17}, \code{small-7b} and
        \code{0.3125} too. A pattern that is too loose finds too much.
  \item \code{findall} returns a list of strings if the pattern has no groups, a list of strings
        if it has one group, and a list of \emph{tuples} if it has two or more.
  \item In \code{sub}, \code{\textbackslash 1}, \code{\textbackslash 2} refer to groups. The
        replacement can also be a \emph{function} that receives the match and returns a string.
\end{itemize}

Other tools you will use daily: \code{re.finditer()} (like \code{findall} but yields
\code{Match} objects with positions), \code{re.fullmatch()} (the whole string must match), and
\code{re.compile(p)} (build the pattern once, reuse it; the compiled object has the same methods).

\begin{pycode}
for mm in re.finditer(r"\d+", "x12 y345"):
    print(mm.group(), mm.start(), mm.end())
print(re.fullmatch(r"\d+", "123"))
print(re.fullmatch(r"\d+", "123a"))
pat = re.compile(r"error", re.IGNORECASE)
print(pat.findall("Error ERROR error"))
\end{pycode}
\begin{rawcode}
12 1 3
345 5 8
<re.Match object; span=(0, 3), match='123'>
None
['Error', 'ERROR', 'error']
\end{rawcode}

\subsection{Patterns}

\subsubsection*{Special characters}

\begin{center}
\begin{tabularx}{\linewidth}{@{}l X@{}}
\toprule
\textbf{Pattern} & \textbf{Matches} \\
\midrule
\code{\textbackslash d} / \code{\textbackslash D} & a digit / a non-digit \\
\code{\textbackslash w} / \code{\textbackslash W} & a letter, digit or underscore / anything else \\
\code{\textbackslash s} / \code{\textbackslash S} & a whitespace character / a non-whitespace character \\
\code{\textbackslash b} / \code{\textbackslash B} & a word boundary / not a word boundary \\
\code{.} & any single character except newline \\
\bottomrule
\end{tabularx}
\end{center}

For \code{str} patterns these classes are Unicode-aware. \code{\textbackslash d} matches digits
from other scripts, and \code{\textbackslash w} matches accented letters. Add the flag
\code{re.ASCII} if you want plain ASCII only.

\begin{pycode}
print(ascii(re.findall(r"\d", "a1\u0663")))
print(ascii(re.findall(r"\d", "a1\u0663", re.ASCII)))
x = "abc-/*\u00ea\u0115"
print(ascii(re.findall(r"\w", x)))
\end{pycode}
\begin{rawcode}
['1', '\u0663']
['1']
['a', 'b', 'c', '\xea', '\u0115']
\end{rawcode}

\subsubsection*{Specifiers}

\begin{center}
\begin{tabularx}{\linewidth}{@{}l X@{}}
\toprule
\textbf{Pattern} & \textbf{Meaning} \\
\midrule
\code{abc} & the literal text \code{abc} \\
\code{(expr)} & group: treat \code{expr} as a unit and capture it \\
\code{a|b} & \code{a} or \code{b} \\
\code{\^{}} \ \ \code{\$} & start / end of the string (anchors) \\
\code{x?} \ \code{x*} \ \code{x+} & zero or one / zero or more / one or more \code{x} \\
\code{x\{m\}} \ \code{x\{m,n\}} & exactly \code{m} / between \code{m} and \code{n} \code{x} \\
\code{x*?} \ \code{x+?} \ \code{x??} & the same, but match as \emph{little} as possible (lazy) \\
\code{[abc]} \ \code{[\^{}abc]} & one of these characters / any character \emph{not} listed \\
\code{x(?=y)} \ \code{x(?!y)} & \code{x} only if followed / not followed by \code{y} \\
\code{(?<=y)x} \ \code{(?<!y)x} & \code{x} only if preceded / not preceded by \code{y} \\
\bottomrule
\end{tabularx}
\end{center}

\begin{pycode}
src = "I wish I may, I wish I might\nHave a dish of fish tonight."
print(re.findall(r"wish|fish", src))
print(re.findall(r"^wish", src), re.findall(r"^I wish", src))
print(re.findall(r"fish$", src), re.findall(r"fish tonight\.$", src))
print(re.findall(r"[wf]ish", src))
print(re.findall(r"[^aeiou\s]+", "hello world"))
print(re.findall(r"I (?=wish)", src))
print(re.findall(r"(?<=I) wish", src))
print(re.findall(r"I (?!wish)", src))
\end{pycode}
\begin{rawcode}
['wish', 'wish', 'fish']
[] ['I wish']
[] ['fish tonight.']
['wish', 'wish', 'fish']
['h', 'll', 'w', 'rld']
['I ', 'I ']
[' wish', ' wish']
['I ', 'I ']
\end{rawcode}

By default \code{\^{}} and \code{\$} look only at the start and end of the \emph{whole} string.
Add \code{re.MULTILINE} to make them work per line, and \code{re.DOTALL} so that \code{.}
also matches a newline.

\begin{pycode}
print(re.findall(r"^\w+", "ab cd\nef gh"))
print(re.findall(r"^\w+", "ab cd\nef gh", re.MULTILINE))
print(re.findall(r"a.b", "a\nb"), re.findall(r"a.b", "a\nb", re.DOTALL))
\end{pycode}
\begin{rawcode}
['ab']
['ab', 'ef']
[] ['a\nb']
\end{rawcode}

\subsubsection*{Greedy versus lazy}

Quantifiers grab as much as they can. That surprises people when they extract text between two
markers.

\begin{pycode}
t = "<b>one</b><b>two</b>"
print(re.findall(r"<b>.*</b>", t))      # greedy: one big match
print(re.findall(r"<b>.*?</b>", t))     # lazy: stops at the first </b>
print(re.findall(r"a{2,3}", "a aa aaa aaaa"))
\end{pycode}
\begin{rawcode}
['<b>one</b><b>two</b>']
['<b>one</b>', '<b>two</b>']
['aa', 'aaa', 'aaa']
\end{rawcode}

\begin{center}
\begin{tikzpicture}[>=Stealth]
  \foreach \c [count=\i from 0] in {<,b,>,o,n,e,<,\slash,b,>,<,b,>,t,w,o,<,\slash,b,>}
    \node[c17cell,minimum width=4.6mm] (s\i) at (\i*0.52,0) {\c};
  \node[font=\sffamily\scriptsize,anchor=east] at (-0.4,0) {source};
  \draw[fill=cred!20,draw=cred,rounded corners=2pt] (-0.25,-0.55) rectangle (9.85,-0.95);
  \node[font=\sffamily\scriptsize,anchor=east] at (-0.4,-0.75) {greedy \code{.*}};
  \node[font=\sffamily\scriptsize] at (5,-0.75) {one match: swallows everything up to the \emph{last} \code{</b>}};
  \draw[fill=cgreen!25,draw=cgreen,rounded corners=2pt] (-0.25,-1.25) rectangle (4.85,-1.65);
  \draw[fill=cgreen!25,draw=cgreen,rounded corners=2pt] (5.0,-1.25) rectangle (9.85,-1.65);
  \node[font=\sffamily\scriptsize,anchor=east] at (-0.4,-1.45) {lazy \code{.*?}};
  \node[font=\sffamily\scriptsize] at (2.3,-1.45) {match 1};
  \node[font=\sffamily\scriptsize] at (7.4,-1.45) {match 2};
\end{tikzpicture}
\end{center}
% MD-BEGIN
% ```text
% source: <b>one</b><b>two</b>
% greedy  <b>.*</b> : [<b>one</b><b>two</b>]            one big match
% lazy    <b>.*?</b>: [<b>one</b>] [<b>two</b>]         two small matches
% ```
% MD-END

\subsubsection*{Always use raw strings for patterns}

In a normal string, \code{"\textbackslash b"} is the backspace character. In a regex,
\code{\textbackslash b} means word boundary. Put an \code{r} before the quote so Python leaves
backslashes alone.

\begin{pycode}
print(len("\b"), len(r"\b"))
print(re.findall("\bcat", "cat"), re.findall(r"\bcat", "cat"))
print(re.findall(r"\bcat\b", "cat concat cat."))
print(re.escape("a.b*c"))
print(re.findall("price: $5", "price: $5"))
print(re.findall(re.escape("price: $5"), "price: $5"))
\end{pycode}
\begin{rawcode}
1 2
[] ['cat']
['cat', 'cat']
a\.b\*c
[]
['price: $5']
\end{rawcode}

\begin{trap}
Forgetting the \code{r} prefix often \emph{works by accident} (\code{"\textbackslash d"} is
left alone by Python) and then fails on \code{\textbackslash b}. Make the \code{r} a habit.
Also escape special characters when you want them literally: use \code{\textbackslash .} for a
dot, or let \code{re.escape()} do it for user-supplied text. The unescaped \code{\$} in the last
example means ``end of string'', so it found nothing.
\end{trap}

\subsection{Specifying match output}

Parentheses make groups. Add \code{?P<name>} inside to name a group; then you read it by
name, or get a ready-made dict.

\begin{pycode}
m = re.search(r"run=(?P<run>\w+) .*loss=(?P<loss>[\d.]+)", line)
print(m.group("run"), m["loss"], m.groupdict())
out = "Answer: 42. Confidence: 0.87"
print(re.search(r"Confidence:\s*([\d.]+)", out).group(1))
\end{pycode}
\begin{rawcode}
a17 0.3125 {'run': 'a17', 'loss': '0.3125'}
0.87
\end{rawcode}

\begin{gamerule}
Use named groups for any pattern with more than one group, and compile it once at module level
with a clear name (\code{LOSS\_RE}). Regexes are write-only if you do not name their parts.
When the text is JSON, parse it with \code{json}, not regex. Use regex for text that has no
better structure.
\end{gamerule}

% ------------------------------------------------------------
\section{Review and what comes next}

You now know that a \code{str} is Unicode code points, that \code{.encode()} and
\code{.decode()} are the only doors between \code{str} and \code{bytes}, and that the \code{re}
module finds, splits and rewrites text by pattern. The next chapter looks at the \code{bytes}
side: binary data.

\begin{practice}
\begin{enumerate}[nosep,leftmargin=*]
  \item Build a string holding the Unicode character with code point \code{U+1F984} using only
        an escape sequence. Print its official name and its length.
  \item Encode that string to UTF-8. How many bytes do you get? Decode them back and compare
        with the original using \code{==}.
  \item Why does \code{"e\textbackslash u0301" == "\textbackslash u00e9"} give \code{False}? Which
        call makes the comparison true?
  \item Write a pattern, with a raw string, that finds every number with a decimal point in the
        text \code{"lr 0.001, wd 0.01, epochs 3"}. What does \code{findall} return?
  \item Given \code{"[A] x [B] y"}, extract only the letters in brackets. First write it greedy so
        it fails, then fix it.
  \item You get \code{b"caf\textbackslash xc3\textbackslash xa9"} from an API. Which line turns it
        into a \code{str}, and what goes wrong if you pass \code{"latin-1"}?
\end{enumerate}
\end{practice}

\begin{achievement}{Text Wrangler}
\begin{itemize}[nosep,leftmargin=*]
  \item Tell \code{str} from \code{bytes} and convert between them with an explicit encoding.
  \item Handle encoding errors on purpose, and recognise mojibake.
  \item Normalize Unicode text before comparing it.
  \item Extract fields from a log line with \code{search}, \code{findall} and named groups.
  \item Write regex patterns as raw strings and avoid the greedy trap.
\end{itemize}
\end{achievement}
